% =====================================================================
% Section: Two-Electron Repulsion on S^3
% For inclusion in Paper 7 (The Dimensionless Vacuum)
% Date: March 2026
% =====================================================================

\section{Two-Electron Repulsion on $S^3$}
\label{sec:vee_s3}

The preceding sections established that the single-electron Coulomb
problem maps exactly onto the unit three-sphere $S^3$ via Fock's
stereographic projection. We now extend this framework to the
two-electron repulsion integral and derive its $S^3$ representation.

\subsection{Momentum-Space Representation}

The direct Coulomb integral (Slater $F^0$) between orbitals $a$ and $b$
in position space is
\begin{equation}
  F^0(a,b) = \int_0^\infty \!\!\int_0^\infty
  |R_a(r_1)|^2 \, |R_b(r_2)|^2 \, \frac{1}{r_>} \,
  r_1^2 \, r_2^2 \, dr_1 \, dr_2 \,,
  \label{eq:slater_f0}
\end{equation}
where $r_> = \max(r_1, r_2)$. For hydrogen-like orbitals with nuclear
charge $Z$, exact values are known:
\begin{align}
  F^0(1s,1s) &= \tfrac{5}{8}Z, \quad
  F^0(1s,2s) = \tfrac{17}{81}Z, \nonumber\\
  F^0(2s,2s) &= \tfrac{77}{512}Z.
  \label{eq:slater_exact}
\end{align}

An equivalent momentum-space representation exists for $s$-orbital pairs.
Defining the density Fourier transform
\begin{equation}
  \tilde{\rho}_a(q)
  = \frac{1}{q} \int_0^\infty |R_a(r)|^2 \, r \, \sin(qr) \, dr \,,
  \label{eq:density_ft}
\end{equation}
the Slater integral becomes
\begin{equation}
  F^0(a,b) = \frac{2}{\pi} \int_0^\infty
  \tilde{\rho}_a(q) \, \tilde{\rho}_b(q) \, dq \,.
  \label{eq:f0_momentum}
\end{equation}

For the hydrogen $1s$ and $2s$ orbitals:
\begin{align}
  \tilde{\rho}_{1s}(q) &= \frac{16Z^4}{(q^2 + 4Z^2)^2} \,,
  \label{eq:rho1s} \\[4pt]
  \tilde{\rho}_{2s}(q) &= \frac{Z^4(Z^4 - 3Z^2 q^2 + 2q^4)}{(Z^2 + q^2)^4} \,.
  \label{eq:rho2s}
\end{align}

\subsection{Fock Projection onto $S^3$}

We apply the Fock stereographic projection with parameter $p_0 = 2Z$,
under which $q = 2Z \tan u$ for $u \in [0, \pi/2)$ and
$dq = 2Z/\cos^2 u \, du$. Define the \emph{projected density}
\begin{equation}
  \Phi_a(t) \equiv \tilde{\rho}_a(2Zt) \,,
  \label{eq:projected_density}
\end{equation}
where $t = \tan u = q/(2Z)$ is the dimensionless momentum transfer.
Under this substitution, Eq.~\eqref{eq:f0_momentum} becomes
\begin{equation}
  \boxed{
    F^0(a,b) = \frac{4Z}{\pi} \int_0^\infty
    \Phi_a(t) \, \Phi_b(t) \, dt
  }
  \label{eq:f0_s3}
\end{equation}
This is the master formula: $F^0$ is a \emph{single} integral over one
$S^3$ coordinate, not a double integral over two points.

\subsubsection{$1s$--$1s$ verification}

Substituting $q = 2Zt$ into Eq.~\eqref{eq:rho1s}:
\begin{equation}
  \Phi_{1s}(t) = \frac{16Z^4}{(4Z^2 t^2 + 4Z^2)^2}
  = \frac{1}{(t^2 + 1)^2} \,.
  \label{eq:phi1s}
\end{equation}
Then
\begin{equation}
  F^0(1s,1s) = \frac{4Z}{\pi} \int_0^\infty
  \frac{dt}{(t^2+1)^4}
  = \frac{4Z}{\pi} \cdot \frac{5\pi}{32}
  = \frac{5Z}{8} \,. \quad\checkmark
\end{equation}

In angular coordinates ($t = \tan u$, $dt = du/\cos^2 u$),
the integrand becomes $\cos^6 u$, showing that the $1s$ density
is \emph{constant} on $S^3$ (up to the conformal Jacobian).

\subsubsection{$1s$--$2s$ verification (critical test)}

Substituting into Eq.~\eqref{eq:rho2s}:
\begin{equation}
  \Phi_{2s}(t) = \frac{1 - 12t^2 + 32t^4}{(4t^2 + 1)^4} \,.
  \label{eq:phi2s}
\end{equation}
This is \emph{not} constant on $S^3$---it has zeros at
$t^2 = (3 \pm 1)/16$, corresponding to nodes in the $2s$ radial
wavefunction. The integral evaluates to
\begin{align}
  F^0(1s,2s) &= \frac{4Z}{\pi} \int_0^\infty
  \frac{1 - 12t^2 + 32t^4}{(t^2+1)^2 (4t^2+1)^4} \, dt \nonumber\\
  &= \frac{4Z}{\pi} \cdot \frac{17\pi}{324}
  = \frac{17Z}{81} \,. \quad\checkmark
\end{align}

\subsubsection{$2s$--$2s$ verification}

\begin{align}
  F^0(2s,2s) &= \frac{4Z}{\pi} \int_0^\infty
  \frac{(1 - 12t^2 + 32t^4)^2}{(4t^2+1)^8} \, dt \nonumber\\
  &= \frac{4Z}{\pi} \cdot \frac{77\pi}{2048}
  = \frac{77Z}{512} \,. \quad\checkmark
\end{align}

All three Slater integrals are reproduced exactly by the $S^3$
density-overlap formula~\eqref{eq:f0_s3}.

\subsection{Structural Result: Node Property vs.\ Edge Property}
\label{sec:node_vs_edge}

The key structural finding is that
Eq.~\eqref{eq:f0_s3} is a \emph{single} integral over one $S^3$
coordinate (the momentum transfer $t$), not a double integral over
two independent $S^3$ points. Each orbital $a$ contributes a density
\emph{function} $\Phi_a(t)$ on $S^3$, and $V_{ee}$ is the weighted
overlap:
\begin{equation}
  V_{ee}(a,b) = \langle \Phi_a | \Phi_b \rangle_{S^3}
  \equiv \frac{4Z}{\pi} \int_0^\infty \Phi_a(t) \, \Phi_b(t) \, dt \,.
\end{equation}

In graph-theoretic language, $V_{ee}$ is a \textbf{node property}
(each node carries a density profile) rather than an \textbf{edge
property} (a pairwise distance between node coordinates).

\paragraph{Failure of the chordal ansatz.}
A naive chordal-distance ansatz $V_{ee} = \kappa_{ee}/d^2_{\text{chord}}$,
where $d^2 = 2(1 - \boldsymbol{\xi}_a \cdot \boldsymbol{\xi}_b)$
is the chordal distance between $S^3$ points assigned to each orbital,
gives (with $\kappa_{ee} = 5Z/2$):
\begin{center}
\begin{tabular}{lccc}
\hline
Pair & Slater $F^0$ & $\kappa/d^2$ & Ratio \\
\hline
$1s$--$1s$ & $5Z/8$ & $5Z/8$ & 1.00 \\
$1s$--$2s$ & $17Z/81$ & $25Z/4$ & 29.8 \\
$2s$--$2s$ & $77Z/512$ & $5Z/8$ & 4.2 \\
\hline
\end{tabular}
\end{center}
The chordal ansatz is calibrated to reproduce $F^0(1s,1s)$ exactly
(because the $1s$ density is constant on $S^3$, making a single point
adequate), but overestimates $F^0(1s,2s)$ by a factor of~30.
The fundamental reason is that the $2s$ density function has nodes
and oscillations on $S^3$ that cannot be captured by a single coordinate
point.

\subsection{Natural Projection Scales}

Each shell has a natural Fock projection scale:
\begin{equation}
  n = 1: \quad p_0 = 2Z \qquad\text{(denominator } q^2 + 4Z^2\text{)} \,,
\end{equation}
\begin{equation}
  n = 2: \quad p_0 = Z \qquad\text{(denominator } q^2 + Z^2\text{)} \,.
\end{equation}
When all orbitals are projected using a \emph{single} $p_0 = 2Z$, the
$n=1$ density becomes the constant function $\Phi_{1s} = 1/(1+t^2)^2$,
while higher-shell densities acquire polynomial modulations. These
modulations encode the radial node structure and are responsible for
the reduced overlap integrals ($17/81 < 5/8$).

\subsection{Implications}
\label{sec:vee_implications}

The density-overlap formula~\eqref{eq:f0_s3} has three implications
for the computational framework:

\begin{enumerate}
  \item \textbf{Exact Slater integrals on $S^3$.}
    The Fock projection converts position-space double radial integrals
    into single rational-function integrals on $S^3$, which can be
    evaluated in closed form via partial fractions.

  \item \textbf{Graph implementation.}
    In the graph Laplacian framework, each node (orbital state) should
    carry a density profile $\Phi_a(t)$, and $V_{ee}$ matrix elements
    should be computed as overlap integrals rather than pairwise edge
    weights.

  \item \textbf{Angular momentum extension.}
    The $1D$ monopole formula~\eqref{eq:f0_momentum} applies to
    $s$-orbital pairs. For $l > 0$, the full $3D$ angular convolution
    is required, involving the multipole expansion of the orbital
    density in spherical harmonics.
\end{enumerate}
